\IfFileExists{stacks-project.cls}{%
\documentclass{stacks-project}
}{%
\documentclass{amsart}
}

% For dealing with references we use the comment environment
\usepackage{verbatim}
\newenvironment{reference}{\comment}{\endcomment}
%\newenvironment{reference}{}{}
\newenvironment{slogan}{\comment}{\endcomment}
\newenvironment{history}{\comment}{\endcomment}

% For commutative diagrams we use Xy-pic
\usepackage[all]{xy}

% We use 2cell for 2-commutative diagrams.
\xyoption{2cell}
\UseAllTwocells

% We use multicol for the list of chapters between chapters
\usepackage{multicol}

% This is generally recommended for better output
\usepackage{lmodern}
\usepackage[T1]{fontenc}

% For cross-file-references
\usepackage{xr-hyper}

% Package for hypertext links:
\usepackage{hyperref}

% For any local file, say "hello.tex" you want to link to please
% use \externaldocument[hello-]{hello}
\externaldocument[introduction-]{introduction}
\externaldocument[conventions-]{conventions}
\externaldocument[sets-]{sets}
\externaldocument[categories-]{categories}
\externaldocument[topology-]{topology}
\externaldocument[sheaves-]{sheaves}
\externaldocument[sites-]{sites}
\externaldocument[stacks-]{stacks}
\externaldocument[fields-]{fields}
\externaldocument[algebra-]{algebra}
\externaldocument[brauer-]{brauer}
\externaldocument[homology-]{homology}
\externaldocument[derived-]{derived}
\externaldocument[simplicial-]{simplicial}
\externaldocument[more-algebra-]{more-algebra}
\externaldocument[smoothing-]{smoothing}
\externaldocument[modules-]{modules}
\externaldocument[sites-modules-]{sites-modules}
\externaldocument[injectives-]{injectives}
\externaldocument[cohomology-]{cohomology}
\externaldocument[sites-cohomology-]{sites-cohomology}
\externaldocument[dga-]{dga}
\externaldocument[dpa-]{dpa}
\externaldocument[sdga-]{sdga}
\externaldocument[hypercovering-]{hypercovering}
\externaldocument[schemes-]{schemes}
\externaldocument[constructions-]{constructions}
\externaldocument[properties-]{properties}
\externaldocument[morphisms-]{morphisms}
\externaldocument[coherent-]{coherent}
\externaldocument[divisors-]{divisors}
\externaldocument[limits-]{limits}
\externaldocument[varieties-]{varieties}
\externaldocument[topologies-]{topologies}
\externaldocument[descent-]{descent}
\externaldocument[perfect-]{perfect}
\externaldocument[more-morphisms-]{more-morphisms}
\externaldocument[flat-]{flat}
\externaldocument[groupoids-]{groupoids}
\externaldocument[more-groupoids-]{more-groupoids}
\externaldocument[etale-]{etale}
\externaldocument[chow-]{chow}
\externaldocument[intersection-]{intersection}
\externaldocument[pic-]{pic}
\externaldocument[weil-]{weil}
\externaldocument[adequate-]{adequate}
\externaldocument[dualizing-]{dualizing}
\externaldocument[duality-]{duality}
\externaldocument[discriminant-]{discriminant}
\externaldocument[derham-]{derham}
\externaldocument[local-cohomology-]{local-cohomology}
\externaldocument[algebraization-]{algebraization}
\externaldocument[curves-]{curves}
\externaldocument[resolve-]{resolve}
\externaldocument[models-]{models}
\externaldocument[functors-]{functors}
\externaldocument[equiv-]{equiv}
\externaldocument[pione-]{pione}
\externaldocument[etale-cohomology-]{etale-cohomology}
\externaldocument[proetale-]{proetale}
\externaldocument[relative-cycles-]{relative-cycles}
\externaldocument[more-etale-]{more-etale}
\externaldocument[trace-]{trace}
\externaldocument[crystalline-]{crystalline}
\externaldocument[spaces-]{spaces}
\externaldocument[spaces-properties-]{spaces-properties}
\externaldocument[spaces-morphisms-]{spaces-morphisms}
\externaldocument[decent-spaces-]{decent-spaces}
\externaldocument[spaces-cohomology-]{spaces-cohomology}
\externaldocument[spaces-limits-]{spaces-limits}
\externaldocument[spaces-divisors-]{spaces-divisors}
\externaldocument[spaces-over-fields-]{spaces-over-fields}
\externaldocument[spaces-topologies-]{spaces-topologies}
\externaldocument[spaces-descent-]{spaces-descent}
\externaldocument[spaces-perfect-]{spaces-perfect}
\externaldocument[spaces-more-morphisms-]{spaces-more-morphisms}
\externaldocument[spaces-flat-]{spaces-flat}
\externaldocument[spaces-groupoids-]{spaces-groupoids}
\externaldocument[spaces-more-groupoids-]{spaces-more-groupoids}
\externaldocument[bootstrap-]{bootstrap}
\externaldocument[spaces-pushouts-]{spaces-pushouts}
\externaldocument[spaces-chow-]{spaces-chow}
\externaldocument[groupoids-quotients-]{groupoids-quotients}
\externaldocument[spaces-more-cohomology-]{spaces-more-cohomology}
\externaldocument[spaces-simplicial-]{spaces-simplicial}
\externaldocument[spaces-duality-]{spaces-duality}
\externaldocument[formal-spaces-]{formal-spaces}
\externaldocument[restricted-]{restricted}
\externaldocument[spaces-resolve-]{spaces-resolve}
\externaldocument[formal-defos-]{formal-defos}
\externaldocument[defos-]{defos}
\externaldocument[cotangent-]{cotangent}
\externaldocument[examples-defos-]{examples-defos}
\externaldocument[algebraic-]{algebraic}
\externaldocument[examples-stacks-]{examples-stacks}
\externaldocument[stacks-sheaves-]{stacks-sheaves}
\externaldocument[criteria-]{criteria}
\externaldocument[artin-]{artin}
\externaldocument[quot-]{quot}
\externaldocument[stacks-properties-]{stacks-properties}
\externaldocument[stacks-morphisms-]{stacks-morphisms}
\externaldocument[stacks-limits-]{stacks-limits}
\externaldocument[stacks-cohomology-]{stacks-cohomology}
\externaldocument[stacks-perfect-]{stacks-perfect}
\externaldocument[stacks-introduction-]{stacks-introduction}
\externaldocument[stacks-more-morphisms-]{stacks-more-morphisms}
\externaldocument[stacks-geometry-]{stacks-geometry}
\externaldocument[moduli-]{moduli}
\externaldocument[moduli-curves-]{moduli-curves}
\externaldocument[examples-]{examples}
\externaldocument[exercises-]{exercises}
\externaldocument[guide-]{guide}
\externaldocument[desirables-]{desirables}
\externaldocument[coding-]{coding}
\externaldocument[obsolete-]{obsolete}
\externaldocument[fdl-]{fdl}
\externaldocument[index-]{index}

% Theorem environments.
%
\theoremstyle{plain}
\newtheorem{theorem}[subsection]{Theorem}
\newtheorem{proposition}[subsection]{Proposition}
\newtheorem{lemma}[subsection]{Lemma}

\theoremstyle{definition}
\newtheorem{definition}[subsection]{Definition}
\newtheorem{example}[subsection]{Example}
\newtheorem{exercise}[subsection]{Exercise}
\newtheorem{situation}[subsection]{Situation}

\theoremstyle{remark}
\newtheorem{remark}[subsection]{Remark}
\newtheorem{remarks}[subsection]{Remarks}

\numberwithin{equation}{subsection}

% Macros
%
\def\lim{\mathop{\mathrm{lim}}\nolimits}
\def\colim{\mathop{\mathrm{colim}}\nolimits}
\def\Spec{\mathop{\mathrm{Spec}}}
\def\Hom{\mathop{\mathrm{Hom}}\nolimits}
\def\Ext{\mathop{\mathrm{Ext}}\nolimits}
\def\SheafHom{\mathop{\mathcal{H}\!\mathit{om}}\nolimits}
\def\SheafExt{\mathop{\mathcal{E}\!\mathit{xt}}\nolimits}
\def\Sch{\mathit{Sch}}
\def\Mor{\mathop{\mathrm{Mor}}\nolimits}
\def\Ob{\mathop{\mathrm{Ob}}\nolimits}
\def\Sh{\mathop{\mathit{Sh}}\nolimits}
\def\NL{\mathop{N\!L}\nolimits}
\def\CH{\mathop{\mathrm{CH}}\nolimits}
\def\proetale{{pro\text{-}\acute{e}tale}}
\def\etale{{\acute{e}tale}}
\def\QCoh{\mathit{QCoh}}
\def\Ker{\mathop{\mathrm{Ker}}}
\def\Im{\mathop{\mathrm{Im}}}
\def\Coker{\mathop{\mathrm{Coker}}}
\def\Coim{\mathop{\mathrm{Coim}}}

% Boxtimes
%
\DeclareMathSymbol{\boxtimes}{\mathbin}{AMSa}{"02}

%
% Macros for moduli stacks/spaces
%
\def\QCohstack{\mathcal{QC}\!\mathit{oh}}
\def\Cohstack{\mathcal{C}\!\mathit{oh}}
\def\Spacesstack{\mathcal{S}\!\mathit{paces}}
\def\Quotfunctor{\mathrm{Quot}}
\def\Hilbfunctor{\mathrm{Hilb}}
\def\Curvesstack{\mathcal{C}\!\mathit{urves}}
\def\Polarizedstack{\mathcal{P}\!\mathit{olarized}}
\def\Complexesstack{\mathcal{C}\!\mathit{omplexes}}
% \Pic is the operator that assigns to X its picard group, usage \Pic(X)
% \Picardstack_{X/B} denotes the Picard stack of X over B
% \Picardfunctor_{X/B} denotes the Picard functor of X over B
\def\Pic{\mathop{\mathrm{Pic}}\nolimits}
\def\Picardstack{\mathcal{P}\!\mathit{ic}}
\def\Picardfunctor{\mathrm{Pic}}
\def\Deformationcategory{\mathcal{D}\!\mathit{ef}}
\usepackage{cleveref}
\crefname{originallemma}{original lemma}{original lemmas}
\Crefname{originallemma}{Original lemma}{Original lemmas}
\makeatletter
\newlabel{lemma-local-source-target-implies}{{04R0}{source}{}{sourcepointer-1}{}}
\newlabel{lemma-local-source-target-implies@cref}{{[originallemma][0][]04R0}{[1][0][]source}{}{}{}}
\makeatother
\title{Seven equivalent formulations of etale localness on the source-and-target}
\author{The Stacks Project Authors}
\date{stacks-project:descent.tex:local-snapshot-20261012; GNU FDL 1.2 or later}
\begin{document}
\maketitle
\noindent\textit{Editorial note.} All delivered mathematical text is the author's own, copied from the pinned Stacks Project chapter file descent.tex; no statement or proof step was rewritten, reordered or removed.\par\medskip
\noindent\textit{Editorial note.} The counted spans carry 2 navigation stubs for the definition and the equivalence lemma whose numbered clauses the proof cites; each application of a clause is written out in the span.\par\medskip
\noindent\textit{Editorial note.} The preamble span points at the chapter's own preamble.tex shipped with the frozen source, used verbatim as staged.\par\medskip
\noindent\textit{Editorial note.} The label lemma-local-source-target-characterize is defined only in descent.tex, and its canonical tag is not among the tags of the published items of this file, so it is not a published duplicate. In particular it is not the same proposition as the published item 634 (tag 04R2, 'Separate etale locality yields source-and-target locality'), which is a different labelled statement of this chapter.\par\medskip
\noindent\textit{Allowed standard prerequisites (editorial).}\begin{itemize}
\item Definition of etale localness on the source-and-target. Same chapter, Definition definition-local-source-target, invoked at proof lines 7502-7503 and 7513-7514. Supplies the definitional equivalence and clause (1).
\item Equivalent formulations of etale localness and their clause numbering. Same chapter, Lemma lemma-local-source-target-implies parts (1), (2), (3) and (4), invoked at proof lines 7504-7507 and 7515-7524. Supplies the covering formulations and the transfers of the property.
\end{itemize}\medskip
\ifdefined\hypertarget\hypertarget{sourcepointer-1}{}\fi
\noindent\textit{Source reference.} Original lemma 04R0: sources/stacks-descent/pinned-original-source.tex, source lines 7399--7399. Pointer outside the counted span, invoked once for context; the load-bearing argument of the counted statement is the author's proof reproduced in full in the proof span.\par\medskip
\begin{definition}
\label{definition-local-source-target}
Let $\mathcal{P}$ be a property of morphisms of schemes.
We say $\mathcal{P}$ is {\it \'etale local on source-and-target} if
\begin{enumerate}
\item (stable under precomposing with \'etale maps)
if $f : X \to Y$ is \'etale and $g : Y \to Z$ has $\mathcal{P}$,
then $g \circ f$ has $\mathcal{P}$,
\item (stable under \'etale base change)
if $f : X \to Y$ has $\mathcal{P}$ and $Y' \to Y$ is \'etale, then
the base change $f' : Y' \times_Y X \to Y'$ has $\mathcal{P}$, and
\item (locality) given a morphism $f : X \to Y$ the following are equivalent
\begin{enumerate}
\item $f$ has $\mathcal{P}$,
\item for every $x \in X$ there exists a commutative diagram
$$
\xymatrix{
U \ar[d]_a \ar[r]_h & V \ar[d]^b \\
X \ar[r]^f & Y
}
$$
with \'etale vertical arrows and $u \in U$ with $a(u) = x$ such that
$h$ has $\mathcal{P}$.
\end{enumerate}
\end{enumerate}
\end{definition}

\begin{lemma}
\label{lemma-local-source-target-characterize}
Let $\mathcal{P}$ be a property of morphisms of schemes which is
\'etale local on source-and-target. Let $f : X \to Y$ be a morphism
of schemes. The following are equivalent:
\begin{enumerate}
\item[(a)] $f$ has property $\mathcal{P}$,
\item[(b)] for every $x \in X$ there exists an \'etale morphism of germs
$a : (U, u) \to (X, x)$, an \'etale morphism $b : V \to Y$, and
a morphism $h : U \to V$ such that $f \circ a = b \circ h$ and
$h$ has $\mathcal{P}$,
\item[(c)]
for any commutative diagram
$$
\xymatrix{
U \ar[d]_a \ar[r]_h & V \ar[d]^b \\
X \ar[r]^f & Y
}
$$
with $a$, $b$ \'etale the morphism $h$ has $\mathcal{P}$,
\item[(d)] for some diagram as in (c)
with $a : U \to X$ surjective $h$ has $\mathcal{P}$,
\item[(e)] there exists an \'etale covering $\{Y_i \to Y\}_{i \in I}$ such
that each base change $Y_i \times_Y X \to Y_i$ has $\mathcal{P}$,
\item[(f)] there exists an \'etale covering $\{X_i \to X\}_{i \in I}$ such
that each composition $X_i \to Y$ has $\mathcal{P}$,
\item[(g)] there exists an \'etale covering $\{Y_i \to Y\}_{i \in I}$ and
for each $i \in I$ an \'etale covering
$\{X_{ij} \to Y_i \times_Y X\}_{j \in J_i}$ such that each morphism
$X_{ij} \to Y_i$ has $\mathcal{P}$.
\end{enumerate}
\end{lemma}

\begin{proof}
The equivalence of (a) and (b) is part of
Definition \ref{definition-local-source-target}.
The equivalence of (a) and (e) is
Lemma \ref{lemma-local-source-target-implies} part (2).
The equivalence of (a) and (f) is
Lemma \ref{lemma-local-source-target-implies} part (1).
As (a) is now equivalent to (e) and (f) it follows that
(a) equivalent to (g).

\medskip\noindent
It is clear that (c) implies (a). If (a) holds, then for any
diagram as in (c) the morphism $f \circ a$ has $\mathcal{P}$ by
Definition \ref{definition-local-source-target} part (1),
whereupon $h$ has $\mathcal{P}$ by
Lemma \ref{lemma-local-source-target-implies} part (4).
Thus (a) and (c) are equivalent. It is clear that (c) implies (d).
To see that (d) implies (a) assume we have a diagram as in (c)
with $a : U \to X$ surjective and $h$ having $\mathcal{P}$.
Then $b \circ h$ has $\mathcal{P}$ by
Lemma \ref{lemma-local-source-target-implies} part (3).
Since $\{a : U \to X\}$ is an \'etale covering we conclude that
$f$ has $\mathcal{P}$ by
Lemma \ref{lemma-local-source-target-implies} part (1).
\end{proof}

\end{document}
